\chapter{十分钟完成第一份文档}\label{ch:quick}
本例只使用内置数字、小数点和几何分数线，无需先导入个人字库。随手册提供 \code{examples/01-first-document.md}，可以导入，也可以直接复制下方源码。

\section{输入与观察}
\begin{enumerate}
\item 打开应用，确认\ui{设置 → 正文}中的内容模式为\ui{Markdown 混排}，Markdown 格式为\ui{普通 Markdown}。
\item 点击左侧源码区，输入或粘贴下列内容。独立公式的起止双美元符号分别占一行。
\item 等待底部显示“预览已更新”。右侧应出现数字正文、两列表格和一个分数；允许出现数学校准提醒。
\end{enumerate}
\par\noindent\begin{minipage}{\linewidth}
\begin{Source}
12.3

| 1 | 2 |
| --- | --- |
| 3.4 | 5.6 |

$$
\frac{1}{2}
$$
\end{Source}
\end{minipage}\par\medskip
\shot{workspace.png}{1}{数字、表格与分数的实际编辑界面。左侧为源码，右侧为手写页面。}

\section{调整一项设置}
将光标放在表格源码中，打开\ui{设置 → 元素 → 表格设置…}。先保留总宽 160 mm、最小行高 8 mm、数字高度 3.8 mm、留白 0.3 mm。确认当前作用范围后修改对齐方式，观察预览，再取消或确认。学习时每次只改一个参数，更容易理解效果。

\section{保存并导出}
\begin{enumerate}
\item 点击顶部\ui{保存}，选择项目名称，保存为 \code{.inkdoc}。下次用\ui{打开 → 打开项目…}继续。
\item 点击\ui{查看问题}。错误需要修复；“数学校准尚未人工复核”提醒表示需要在正式书写前检查纸面效果。
\item 点击顶部\ui{导出 SVG…}，选择一个尚不存在的输出目录。多页文档会按页输出。
\item 核对导出目录中的 \code{page-001.svg} 和说明清单。它们是机器输出；项目文件仍用于后续修改。
\end{enumerate}

\note{不要将带中文标题的示例作为“无需字库”的首次测试。中文能出现在输入框中，但预览仍需要相应手写轨迹。首次练习成功后，再按第\ref{ch:resources}章导入个人资源。}
